% Chapter 8 — the market's clock.
% Sections live in this directory; figures and macros regenerate via
% scripts/ch08/gen_figures.py from data/ (the frozen snapshot's ATM
% ladders -- no smile is refitted) plus deterministic synthetic
% constructions.  Notation additions are in NOTATION.md.
%
% LESSON GOALS (pedagogy contract).  After this chapter the reader can:
%   1. Explain that an implied volatility is a fraction -- a
%      price-determined total variance divided by a chosen clock -- and
%      prove that relabeling the clock changes readings but never prices,
%      butterfly validity, or calendar order.
%   2. Compute the day-weighted variance clock by hand from an event
%      calendar, and the crush factor sqrt(t/tau); explain the overnight
%      vol crush as a scheduled reading correction, not a price move.
%   3. Interpolate a term structure linearly in tau and say exactly what
%      the calendar-clock interpolation invents between quotes.
%   4. Set up the inverse problem of reading a clock off prices -- the
%      peak rule: an interval that accrues materially faster than the
%      higher of its neighbours carries an event of exactly the size
%      that brings it down to that neighbour -- and state its limits:
%      dates inside an interval are unidentified, and an event below
%      max(half a day, 3% of its interval's variance) is invisible.
%   5. Diagnose a real board: locate the earnings interval in the frozen
%      NVDA ladder, explain why the index board yields nothing,
%      and know when to stop believing the rule.
% Every section serves one of these five goals.
\renewcommand{\figdir}{figures/ch08}
\chapter[The market's clock]{The Market's Clock}
\label{ch:clock}
\chaptermark{The market's clock}

\chaptersynopsis{Since Chapter~2 every smile in this book has carried a
maturity variable $\tau$ called variance time, and every chapter so far
has quietly set it equal to the calendar.  This chapter separates
them.  A scheduled event --- an earnings report, a macro print
--- packs several ordinary days' variance into one calendar day, and a
volatility read against the calendar then misbehaves in a recognizable
way: the frozen NVDA board opens the chapter with an 18-day expiry
quoted below both of its neighbours --- more than four volatility
points below the shorter one.
Nothing is wrong with those prices.  The lesson is that implied
volatility is a fraction --- a price-determined total variance divided
by a clock --- and the clock is a modeling choice.  A short computation
with independent normal days shows that the law at expiry knows only
the variance budget, not its schedule; a classical theorem extends this
to every continuous martingale.  The chapter then builds the
day-weighted event clock, proves that relabeling time preserves prices,
butterfly validity and calendar order while volatility readings scale
by exactly $\sqrt{t/\tau}$, and works through an earnings week on paper
--- the famous overnight volatility crush appears as a scheduled
correction of a reading, with no price moving at all.  The clock then
serves twice: dense term-structure curves must be interpolated
on it, or a known event smears into phantom volatility at unquoted
maturities; and the clock itself can be read off the board's own
prices, by one rule: an interval that accrues materially faster than
the higher of its neighbours carries an event, of exactly the size
that brings it down to that neighbour.  That inverse problem has
measured limits --- a floor set by the listing's spacing, dates that
no data can pin inside an interval --- and the chapter closes by
running it on the frozen boards: the rule hands NVDA's earnings
interval four extra days, hands the index board nothing, and, if
allowed candidates everywhere, reads an ordinary quarter as a peak
too.  The calendar it installs is an estimate to review, not an
oracle.}

\section{Three expiries and a broken ruler}
\label{sec:clkboard}

Part~II has been rebuilding the observation, one input at a time.
Chapter~6 fitted the forward and the discount; Chapter~7 removed the
early-exercise right and handed every quote back as a European price.
One input has ridden along unexamined since Chapter~2: the maturity
variable $\tau$ that sits under every total variance $w=\sigma^2\tau$.
The book has been calling it \emph{variance time} from the start ---
while quietly setting it equal to the calendar, since no chapter yet
had a reason to distinguish them.  Chapter~7 closed by naming the
debt: ``two clocks this book has so far been allowed to conflate,
because its frozen data happens to agree on both.''  This chapter is
about the days when they disagree.

Start with real prices, because the disagreement is sitting in the
book's own frozen data.  The snapshot behind every chapter --- the
Monday evening of 2026-08-03 --- carries two names, and so far the
chapters have worked mostly on SPY.  Turn instead to NVDA, the
snapshot's single-stock name, and read its at-the-money implied
volatility across the eight quoted expiries --- the readings computed
as $\sqrt{w/t}$, total variance over calendar time, which this
chapter will call the \emph{calendar reading} and write
$\sigma_{\rm cal}$.  (At the money means at $k=0$, the forward; the
precise read-off rule is one line in
appendix~\ref{app:clkprotocol}, which also lists the eight expiry
dates.  As before, one \emph{vol bp} is $10^{-4}$ of absolute
volatility; one \emph{variance bp} is $10^{-4}$ of total variance.)
Three of the eight readings make the puzzle by themselves:
\[
\underbrace{\MacClkBoardVolFourDPct\%}_{\text{4 days}}
\qquad
\underbrace{\MacClkBoardVolEighteenDPct\%}_{\text{18 days}}
\qquad
\underbrace{\MacClkBoardVolFortySixDPct\%}_{\text{46 days}} .
\]
The same underlier, the same afternoon, three expiries two and four
weeks apart --- and the middle one is quoted
\MacClkBoardDipBp\ vol bp below the shortest.  Is NVDA's volatility
falling or rising with maturity?  The answer is that the
question is broken, and the goal of this section is to see why.

Ask the question in the right units.  A volatility is a rate; before
trusting a rate, look at the quantity it is the rate \emph{of}.  The
quantity is total implied variance $w=\sigma^2 t$ --- for now we write
$t$, the calendar year fraction of Chapter~6, because the calendar is
the only clock we have so far.  The three expiries carry
\[
w(\text{4d})=\MacClkBoardWFourBp,\qquad
w(\text{18d})=\MacClkBoardWEighteenBp,\qquad
w(\text{46d})=\MacClkBoardWFortySixBp
\quad\text{variance bp.}
\]
Now do the division a first course would do: variance gained per
calendar day, interval by interval.  Between the 4-day and the 18-day
expiries, \MacClkBoardLullDays\ calendar days pass, so
\[
\frac{\MacClkBoardWEighteenBp-\MacClkBoardWFourBp}
     {\MacClkBoardLullDays}
=\MacClkBoardLullPerDay\ \text{variance bp per day;}
\]
between the 18-day and the 46-day expiries, \MacClkBoardHotDays\ days
pass, and
\[
\frac{\MacClkBoardWFortySixBp-\MacClkBoardWEighteenBp}
     {\MacClkBoardHotDays}
=\MacClkBoardHotPerDay\ \text{variance bp per day.}
\]
The days are not alike.  A day in the second window carries
\MacClkBoardHotOverLull\ times the variance of a day in the first.
And the second window is not anonymous.  The 18-day expiry is
August~21; the 46-day expiry is September~18; and NVDA reports
earnings once a quarter, on a schedule published well in advance,
with the next report due in late August --- after the first of those
dates, before the second.  The market is quoting exactly what it
believes: the report's day will move the stock like several ordinary
days, every expiry that \emph{contains} that day inherits its
variance, and the expiry that \emph{dodges} it --- August~21 --- gets
to be cheap.

\Cref{fig:clkboard} shows the whole board.  In panel~(a), the orange
curve is the calendar reading $\sigma_{\rm cal}$ of all eight
expiries against calendar days to expiry, on a logarithmic axis.
Notice the shape around the annotated dip: the two shortest expiries
(2 and 4 days, both before August~21) read above $43\%$, the 18-day
expiry reads \MacClkBoardVolEighteenDPct\%, and the 46-day expiry
recovers most of the drop, at \MacClkBoardVolFortySixDPct\%.
Further out, the curve wanders rather than trends --- the highest
reading on the whole board belongs to the 228-day expiry.  Read as
``the market's forecast of NVDA's volatility as a function of
horizon,'' the curve is erratic.  Panel~(b) draws the same eight
quotes as what they actually pin down: accrual rates
$\Delta w_i/\Delta t_i$, as horizontal steps --- one step per
interval, the first interval running from the valuation date (where
accrued variance is zero) to the first expiry, each later one
between consecutive expiries.  (This per-interval rate is called
\emph{forward variance}, and \cref{sec:clkinverse} gives it a symbol
and a starring role.)  The two computations above are the third and
fourth steps: the low step is the pre-report lull, and the shaded
interval --- the one containing the scheduled report --- runs hot.
Panel~(b) is panel~(a) with the division done interval by interval:
nothing here is erratic, it is one hot window sitting among cooler
ones.

\begin{figure}[htbp]
\centering
\paperfig{\textwidth}{fig_clk_board.pdf}
\caption{The opening puzzle (NVDA, frozen snapshot of 2026-08-03).
(a)~At-the-money implied volatility against calendar days to expiry:
the 18-day expiry --- the one that expires \emph{before} the scheduled
late-August earnings report --- is quoted \MacClkBoardDipBp\ vol bp
below the 4-day reading.  (b)~The same board as accrual rates: total
variance gained per calendar year in each interval between consecutive
expiries.  The interval that contains the report (shaded) accrues
\MacClkBoardHotOverLull$\times$ the variance per day of the lull
interval before it.}
\label{fig:clkboard}
\end{figure}

Why call this a \emph{broken ruler} rather than an interesting market
fact?  Because of what it does to every tool the book has built.  Each
single-expiry fit of Part~I is untouched --- a smile model fits its own
expiry's quotes and never asks how fast variance accrued on the way ---
but everything that couples expiries inherits the kink.  A
term-structure view drawn through these readings shows volatility
falling four points in two weeks and then climbing back, which no
trader believes is a forecast.  An interpolation between the quoted
expiries must invent readings for unlisted maturities, and
\cref{sec:clkinterp} shows the calendar-clock rule inventing them
badly.  The local volatility field of Chapter~4 consumes
$\partial_\tau w$ directly --- its numerator \emph{is} the accrual
rate of panel~(b).  Even the calendar-order condition of Chapter~2
(total variance non-decreasing in maturity at fixed $k$) offers no
comfort: it holds on this board, and the board is still a mess to
read, because order says nothing about the accrual being bunched
into one window.

There are two ways to respond.  One is to enrich the model: add a jump
of random size at the report date, price options on the mixture, and
price with jump models from here on.  That road is real
but expensive, and this book does not take it (\cref{sec:clkfrozen}
says where it lives in the literature).  The other response is older
and much cheaper: keep the model, change the \emph{clock}.  If an
earnings day is worth several ordinary days of variance, then measure
time in units in which it \emph{is} several days.  On that clock,
variance accrues at a steady rate again, and the constant-rate
diffusion picture --- the one behind every formula in this book ---
fits without a fight.  The next section shows that this is not a
trick: it is how the mathematics of martingales says time should be
measured in the first place.

\section{The walk keeps its own clock}
\label{sec:clkwalk}

We know what the board says: some days carry more variance than
others.  This section shows that the mathematics has a natural home
for that fact.  The claim to be earned, first in a discrete
computation a first course can check and then as a named theorem, is
this: \emph{the law of the underlying at expiry depends on the days
only through their total variance --- the budget --- and not on how
that budget is scheduled.}  Once that is true, ``which clock to use''
becomes a genuine choice, because the prices cannot object.

Do the discrete case honestly.  Let the log return of the underlying
over trading day $i$ be a random shock $\varepsilon_i$, normally
distributed with mean zero and variance $\Delta w_i$ --- day $i$'s
contribution to the total.  (Mean zero is a deliberate
simplification: it drops the $-\tfrac12$ drift correction Chapter~5
derived, which shifts the mean and not the variance; the correction
reappears, intact, in the continuous computation below.)  Let the
days be independent: today's surprise does not read yesterday's.  After $n$ days the cumulative log return is the sum
$\varepsilon_1+\cdots+\varepsilon_n$, and two first-course facts
finish the argument.  A sum of independent normals is normal.  And
its variance is the sum of the variances:
\[
\operatorname{Var}\!\big(\varepsilon_1+\cdots+\varepsilon_n\big)
=\Delta w_1+\Delta w_2+\cdots+\Delta w_n .
\]
The right-hand side is a \emph{sum}, and sums do not care about order
or arrangement.  Put the big day first, last, or in the middle; split
it into two medium days; the terminal law is the same normal with the
same variance.  The law at expiry knows the budget, not the schedule.

Make it concrete with the numbers this section's figure uses:
\MacClkWalkDays\ trading days, each ordinary day's shock worth
$\MacClkWalkDailyPct\%$ of standard deviation --- that is,
$(\MacClkWalkDailyPct\%)^2$ of variance --- and one scheduled event on
day \MacClkWalkEventDay\ carrying \MacClkWalkEventUnits\ ordinary
days' variance in one calendar day.  Take one ordinary day's variance
as the unit.  The \MacClkWalkDays\ days are 29 ordinary ones plus the
event day; the ordinary days contribute one unit each and the event
day contributes \MacClkWalkEventUnits, so the terminal standard
deviation is
\[
\MacClkWalkDailyPct\%\times
\sqrt{\underbrace{29\times1}_{\text{ordinary days}}
+\underbrace{1\times\MacClkWalkEventUnits}_{\text{the event day}}}
=\MacClkWalkDailyPct\%\times\sqrt{34},
\]
and it would be $\MacClkWalkDailyPct\%\times\sqrt{34}$ with the event
on day 3, day 20, or day 30.  An option struck on the terminal level
--- any option, since Chapter~2 taught that a full strike ladder pins
the terminal law and nothing more --- cannot tell those calendars
apart.

But something \emph{can} tell them apart: any reading that divides by
elapsed calendar time.  \Cref{fig:clkwalk} shows the same simulated
walk twice.  In panel~(a), the horizontal axis is the trading day.
The blue band with grey edges is the $\pm1$-standard-deviation
envelope, $\pm\MacClkWalkDailyPct\%\sqrt{\text{variance so far}}$:
the statistical corridor the walk lives in.  Follow the band's upper edge
with your eye: it grows like a square root, then takes a visible
step up at the dotted line --- day \MacClkWalkEventDay, where one
day adds five days' variance --- and resumes its square-root growth
from a higher base.  The kink is the calendar clock failing to keep
pace with the walk.  In panel~(b), the \emph{same} paths are drawn
against a different horizontal axis: accumulated variance, measured
in \emph{variance days} --- ordinary days count one, the event day
counts \MacClkWalkEventUnits.  The event day is now simply a wider
slot on the axis (the orange strip), the envelope is one smooth
square root from end to end, and nothing about the walk marks the
event at all.  Same randomness, same terminal law; the kink lived in
the ruler.

\begin{figure}[htbp]
\centering
\paperfig{\textwidth}{fig_clk_walk.pdf}
\caption{One walk, two rulers (\MacClkWalkPaths\ simulated paths, one
highlighted; construction and seed in
appendix~\ref{app:clkprotocol}).  (a)~Against the trading-day index,
the $\pm1$-sd envelope (grey edges, blue band) kinks at the event
day: day \MacClkWalkEventDay\ carries \MacClkWalkEventUnits\ days'
variance.  (b)~Against accumulated variance the event day is a wider
slot on the axis (the orange strip), and the envelope is one smooth
square root.  The walk itself never noticed the event's position;
only the ruler did.}
\label{fig:clkwalk}
\end{figure}

The continuous version of this bookkeeping is a classical theorem.
Chapter~5 wrote the underlying's motion as a diffusion,
$\dd Y_s=\sqrt{v_s}\,Y_s\,\dd W_s$ --- $Y$ the forward-normalized
level, $W$ a Brownian motion (the continuous limit of an unbiased
random walk), $v_s$ the instantaneous variance, the squared
volatility the process runs at instant $s$ (\cref{sec:vsnumber}).
There the time variable was variance time itself; run it now on the
\emph{calendar}, where $v_s$ is free to schedule itself unevenly ---
quiet Tuesdays, a violent earnings evening.  Chapter~5's It\^o
computation \eqref{eq:vsito} splits the log return into a drift and a
martingale part --- a \emph{martingale} being, as before, a process
with no drift, whose expected future value is its current value:
\[
\log Y_t
=\underbrace{\int_0^{t}\sqrt{v_s}\,\dd W_s}_{\textstyle M_t,
\ \text{the martingale part}}
\;-\;\frac12\underbrace{\int_0^{t}v_s\,\dd s}_{\textstyle
\mathcal{C}(t),\ \text{accumulated variance}} .
\]
The second integral, written $\mathcal{C}(t)$, is the continuous
analogue of ``variance days'': the running total of instantaneous
variance, growing fast through violent stretches and slowly through
quiet ones.  It is the clock the market itself keeps.  For a general
continuous martingale --- one not handed to us as an integral against
a Brownian motion --- the accumulated variance is defined directly
from the path, as the small-step limit of the sum of squared
increments; for our $M_t=\int_0^t\sqrt{v_s}\,\dd W_s$ that limit is
exactly the $\mathcal{C}(t)=\int_0^t v_s\,\dd s$ above.  With the
term defined, the theorem says the martingale part is nothing but a
Brownian motion read on that clock:

\begin{theorem}[Dambis; Dubins--Schwarz
\citep{Dambis1965,DubinsSchwarz1965}]
\label{thm:clkdds}
Let $M$ be a continuous martingale with $M_0=0$ whose accumulated
variance $\mathcal{C}$ grows without bound.  Then there is a standard
Brownian motion $Z$ such that
\[
M_t \;=\; Z_{\mathcal{C}(t)}\qquad\text{for all }t .
\]
\end{theorem}

Two remarks keep the statement honest at this book's level.  The
growth condition just ensures the relabeled clock runs on forever, so
that $Z$ is a full Brownian motion; it is harmless here, where the
market keeps moving past any one expiry.  And the delicate step ---
reading a Brownian motion at the \emph{random} time $\mathcal{C}(t)$
and getting the original martingale back --- is precisely what the
proof establishes, by a careful change-of-time argument in the
standard reference \citep[Ch.~V]{RevuzYor1999}.  We use the
statement, not the proof.
What it buys us is the discrete computation again, now for every
continuous martingale at once: running the motion to calendar date
$t$, the log return is
$\log Y_t=Z_{\mathcal{C}(t)}-\tfrac12\mathcal{C}(t)$, so when the
schedule $v_s$ is deterministic the terminal law is normal with
variance $\mathcal{C}(t)$ --- the budget --- and mean
$-\tfrac12\mathcal{C}(t)$, no matter how the accrual was arranged
along the way.  The market keeps one clock, $\mathcal{C}$; the
calendar is a second clock that we brought.

Now say what this means for implied volatility, in one careful
paragraph.  A quoted option price, pushed through the Black formula
of Chapter~2, delivers a \emph{total variance} $w(k)$ --- and that
inversion never asks about time at all.  This is worth pausing on,
because a first course writes Black--Scholes with $\sigma$ and $T$ as
separate inputs.  Chapter~2's normalized form \eqref{eq:black} shows
they never act separately: the formula is $B(k,w)$ with
$d_\pm=-k/\sqrt{w}\pm\sqrt{w}/2$, and maturity enters only through
the product $w=\sigma^2\tau$.  Everything after the inversion is
division: quoting the same $w$ as a volatility requires choosing a
clock and computing $\sqrt{w/\text{clock}}$.  Divide by the calendar
$t$ and we write
\[
\sigma_{\rm cal}=\sqrt{w/t}\,;
\]
divide by a clock $\tau$ built to follow the market's own accrual
schedule and we get $\sigma=\sqrt{w/\tau}$ --- which is precisely the
$\sigma$ this book has used since Chapter~2, where $w=\sigma^2\tau$
was the founding convention.  Until today the two clocks agreed,
because no event calendar was in force; the NVDA board is where they
part.  The dip of \cref{sec:clkboard} is now diagnosed: the 18-day
expiry dodges a scheduled burst of $\mathcal{C}$, its $w$ is
accordingly modest, and dividing that modest $w$ by an
\emph{undiscounted} 18 days of calendar produces a number that looks
like a volatility collapse.  The prices are consistent; the reading
$\sigma_{\rm cal}$ is the artifact.

One honest boundary before building the clock.  The theorem licenses
a \emph{relabeling} of time; it does not predict $\mathcal{C}$.  The
market's clock is random --- nobody knows Tuesday's variance in
advance --- and what a fitting desk can carry is a deterministic
\emph{forecast of its schedule}: ordinary days advance the clock at
rate one, scheduled event days advance it by more.  Building that
forecast, as a small object a desk can quote and edit, is the next
section's job.

\section{The day-weighted clock}
\label{sec:clkclock}

We now know the market keeps its own clock, and that a desk can only
carry a forecast of its schedule.  This section builds that forecast
in the simplest currency a trading desk speaks: \emph{days}.

Each scheduled event --- an earnings report, a central-bank decision,
a court ruling --- is recorded as a pair $(t_e, N_e)$: its date
$t_e$, a calendar year fraction like any expiry's, and its size
$N_e$, the number of \emph{extra} equivalent days the event is worth.
Mind the word extra: an earnings day with $N_e=4$ counts as five
ordinary days of variance --- its own one, plus four extra.  The
walk's event day in \cref{sec:clkwalk}, five variance units in one
calendar day, is in this currency an event of size $N_e=4$; the
five-unit-wide strip of \cref{fig:clkwalk}(b) and the four-day jump
about to appear in \cref{fig:clkclock}(a) describe the same event in
the two bookkeepings.  Given a calendar of such events, the variance
time of maturity $t$ is defined by counting days:
\begin{equation}
\boxed{\;\tau_{\rm days}(t)\;=\;365\,t
\;+\!\!\sum_{t_e\le t}\!N_e,
\qquad
\tau(t)\;=\;\frac{\tau_{\rm days}(t)}{365}.\;}
\label{eq:clkclock}
\end{equation}
The sum runs over events at or before $t$: an expiry pays for exactly
the events it contains.  With an empty calendar, $\tau(t)=t$ --- the
convention every previous chapter ran under, now visible as the
special case.

Compute one by hand.  Take a 14-day expiry and a single event worth
$N_e=\MacClkClockExtraDays$ extra days, scheduled
\MacClkClockEventDay\ days out.  The event is inside the expiry, so
\[
\tau_{\rm days}=14+\MacClkClockExtraDays=18,
\qquad
\tau=\frac{18}{365}=0.0493,
\]
against a calendar $t=14/365=0.0384$.  Eighteen variance days will
elapse in fourteen calendar days.  For an expiry of 9 days --- one
that \emph{dodges} the event --- the sum is empty and
$\tau_{\rm days}=9$: nothing changes.  That is the whole mechanism;
everything else in the chapter is consequences.

\Cref{fig:clkclock}(a) draws \eqref{eq:clkclock} for this calendar.
The blue line is $\tau_{\rm days}$ against calendar days: it tracks
the dashed calendar diagonal exactly until day \MacClkClockEventDay,
jumps by \MacClkClockExtraDays\ days there, and runs \emph{parallel}
to the diagonal afterwards --- an event shifts the clock once and is
never paid for again.  (The dotted variant is the normalized clock,
explained at the end of this section; at this zoom it hugs the plain
one.)  Panel~(b) is read after the proposition below.

What may a change of clock do, and what may it never do?  The answer
deserves to be a proposition, because it is the safety case for
everything that follows.  Recall two terms from Chapter~2, one line
each: a smile is \emph{butterfly-valid} when its implied density is
nonnegative --- the Durrleman check $g_D\ge0$ of \eqref{eq:durrleman}
--- and two expiries are in \emph{calendar order} when total variance
does not decrease with maturity at fixed log-moneyness
(\cref{sec:calendar}).

\begin{proposition}[Relabeling time: invariance and covariance]
\label{prop:clkgauge}
Let $\tau(\cdot)$ be any strictly increasing relabeling of maturity
with $\tau(0)=0$ --- in particular the clock \eqref{eq:clkclock}.
Then:
\begin{enumerate}[label=(\roman*)]
\item option prices, and with them every total variance $w(k)$ and
the terminal law of each expiry, are unchanged;
\item butterfly validity of each smile and calendar order across
smiles are unchanged;
\item the volatility reading is \emph{covariant} --- it does not stay
fixed but transforms with the clock, by a known factor:
$\sigma=\sqrt{w/\tau}=\sigma_{\rm cal}\sqrt{t/\tau}$;
\item the accrual rate between consecutive expiries is covariant in
the same sense, by the interval's clock ratio:
$\Delta w/\Delta\tau
=(\Delta w/\Delta t)\cdot(\Delta t/\Delta\tau)$.
\end{enumerate}
\end{proposition}

\begin{proof}
A relabeling attaches new time coordinates to the same expiries; no
random variable, no expectation, and no price changes, which is (i).
Butterfly validity is a property of one expiry's law, unchanged by
(i); calendar order compares the same pairs $w(k,T_1)\le w(k,T_2)$ in
the same order because the relabeling is strictly increasing, which
is (ii).  For (iii) and (iv), divide the same invariant quantities
($w$, $\Delta w$) by the relabeled denominators and cancel:
$\sqrt{w/\tau}=\sqrt{w/t}\cdot\sqrt{t/\tau}$, and likewise for the
ratios of increments.
\end{proof}

Everything a desk can trade sits in the invariant column; everything
the clock touches is a \emph{reading}.  Part~(iv) looks like bare
cancellation, and that is exactly its content: the increments
$\Delta w$ carry all the information, and a clock change re-divides
them without adding any --- a fact \cref{sec:clkinverse} will lean
on.  In particular, installing an event can never move a price and
never create an arbitrage --- not because an implementation is
careful, but because the clock enters the mathematics in exactly one
place: as a denominator.

Part (iii) is worth turning around, because it is the formula desks
feel.  Fix an expiry $t$ and insert one event of $N_e$ extra days
before it.  Total variance $w$ is pinned by the price, so the clock
reading drops by exactly the ratio of the two clocks:
\begin{equation}
\sigma
=\sigma_{\rm cal}\sqrt{\frac{t}{\tau}}
=\sigma_{\rm cal}\sqrt{\frac{365\,t}{365\,t+N_e}} .
\label{eq:clkcrush}
\end{equation}
Keep the direction straight, once, on the hand example.  The clock
reading is the \emph{smaller} number: \eqref{eq:clkcrush} multiplies
the calendar reading by $\sqrt{14/18}$, taking $34.0\%$ down to
$30.0\%$.  Turned over, the calendar reading is the clock reading
\emph{lifted} by
$\sqrt{18/14}=\MacClkClockRatioTwoWeeks$: the same pair of numbers,
the same price.  How big is the effect in general?  For an event small
against the maturity ($N_e\ll365\,t$), expand the square root with
the binomial approximation $\sqrt{1+x}\approx1+\tfrac{x}{2}$, whose
error is of order $x^2$:
\[
\frac{\sigma_{\rm cal}}{\sigma}
=\sqrt{1+\frac{N_e}{365\,t}}
\;\approx\;1+\frac{N_e}{2\cdot365\,t}.
\]
The lift of the calendar reading scales like $N_e$ over \emph{twice
the calendar days to expiry}: four extra days lift a one-year reading
by about half a percent of itself ($4/730$), and a two-week reading
by roughly a seventh ($4/28$; the exact lift is
$\MacClkClockRatioTwoWeeks-1$, a touch less).  Events are a
short-maturity phenomenon --- exactly where desks watch implied
volatility most closely.

Now read \cref{fig:clkclock}(b), which plots the lift direction: the
blue curve is $\sigma_{\rm cal}/\sigma=\sqrt{\tau/t}$ across
expiries, for the same one-event calendar.  Expiries before day
\MacClkClockEventDay\ dodge the event: the ratio is exactly one.  The
expiry \emph{on} the event date pays the full toll,
$\sqrt{14/10}=\MacClkClockPeakRatio$ --- the annotated peak.  Past it
the ratio decays toward one as the event becomes a smaller and
smaller fraction of the maturity.  The dashed curve is the
small-event approximation just derived: generous at the peak, where
$N_e$ is anything but small against $365\,t$, still visibly above
the exact curve out to about a month, and indistinguishable beyond.

\begin{figure}[htbp]
\centering
\paperfig{\textwidth}{fig_clk_clock.pdf}
\caption{The day-weighted clock, one event of
\MacClkClockExtraDays\ extra days at day \MacClkClockEventDay.
(a)~$\tau_{\rm days}(t)$: the calendar diagonal, a jump of
\MacClkClockExtraDays\ days at the event, parallel growth after; the
normalized variant (dotted) rescales the whole line by
\MacClkClockNormFactor.  (b)~The reading ratio
$\sigma_{\rm cal}/\sigma=\sqrt{\tau/t}$ across expiries: one before
the event, a peak of \MacClkClockPeakRatio\ at the event-day expiry,
decay like $1+N_e/(2\cdot365\,t)$ (dashed) beyond it.}
\label{fig:clkclock}
\end{figure}

One bookkeeping choice completes the clock.  As defined,
\eqref{eq:clkclock} makes a year with events \emph{longer} than 365
variance days: the events add variance on top of an unchanged
baseline.  Some desks prefer the opposite convention: the year's
total budget is what it is, and events say \emph{when} it accrues,
not that there is more of it.  The normalized clock implements that
by rescaling all days --- ordinary and event alike --- by one factor
that restores the annual budget:
\begin{equation}
\tau_{\rm days}(t)\;\longmapsto\;
\tau_{\rm days}(t)\cdot
\frac{365}{365+\sum_{t_e\le1}N_e},
\label{eq:clknorm}
\end{equation}
where the sum in the denominator collects the events of the first
year ($t_e\le1$).  For the worked calendar the factor is
$365/369=\MacClkClockNormFactor$ --- the dotted line of
panel~(a).  Under \eqref{eq:clknorm} the one-year variance time, and
with it the one-year volatility reading, exactly match a no-event
year.  Sub-year readings shift slightly: the rescaling makes every
day worth a little less than before, so a maturity that dodges all
events now counts marginally fewer variance days, while the events'
\emph{relative} weight within the year is exactly what it was.
Neither convention is more correct:
un-normalized treats events as extra variance against a
one-per-day baseline, normalized treats the annual budget as known
and the events as its scheduled lumps.  What matters --- the book's
recurring refrain --- is that the choice is stated.  The chapter's
contract table records it, and the whole rest of the chapter runs
un-normalized.

We now own the clock: a desk-sized object --- a list of dated day
counts --- that turns the erratic readings of \cref{sec:clkboard}
into steady ones without touching a single price.  The next section
watches it work through the week every option trader knows best.

\section{An earnings week on paper}
\label{sec:clkcrush}

The clock is built; watch it explain the most famous pattern in
listed options.  Every trader knows the shape: implied volatility on
a name climbs into its earnings date, day after day, and collapses
the morning after --- the \emph{vol crush}.  The usual telling makes
it sound like a repricing, as if the market changed its mind
overnight.  This section constructs an earnings week where the market
never changes its mind about anything, and the whole pattern appears
anyway --- in the readings.

The construction could not be plainer.  A synthetic name runs at a
flat clock volatility of $30\%$ --- on the variance clock, every
expiry, every day, $\sigma=30\%$ exactly --- with one scheduled event
worth \MacClkClockExtraDays\ extra days at day \MacClkClockEventDay.
All prices follow from Chapter~2's formulas with
$w(t)=(0.30)^2\,\tau(t)$ and the clock \eqref{eq:clkclock}; there is
no other input.  Whatever patterns appear below are therefore not
market opinions: they are arithmetic.

\Cref{fig:clkcrush}(a) is the view across expiries, priced at the
start of the week.  The blue line is the clock reading: flat at
$30\%$ by construction.  The orange curve is the calendar reading
$\sigma_{\rm cal}(t)=\sigma\sqrt{\tau(t)/t}$ of the \emph{same
prices}, with dots at the weekly listed expiries.  Expiries before
day \MacClkClockEventDay\ dodge the event and read exactly $30\%$ ---
the annotated 7-day dot.  The expiry on the event day pays the whole
toll at its shortest possible maturity and reads
\MacClkCrushTermPeakPct\% --- the hump's peak, the ratio
\MacClkClockPeakRatio\ of \cref{fig:clkclock}(b) applied to $30\%$.
(The peak falls between listed expiries; the weekly dots at days 7
and 14 straddle it.)  Beyond it the hump decays as the event thins
out over more calendar: by eight weeks the reading is within about a
point of the clock's.  This
is the shape the NVDA board of \cref{sec:clkboard} showed in the
wild, reproduced here from a two-line construction --- with the dip
now explained as the dodging expiries' side of the hump.

\begin{figure}[htbp]
\centering
\paperfig{\textwidth}{fig_clk_crush.pdf}
\caption{An earnings week on paper: flat $30\%$ clock volatility, one
event of \MacClkClockExtraDays\ extra days at day
\MacClkClockEventDay.  (a)~Across expiries: the calendar reading
(orange, weekly dots) is flat at $30\%$ where expiries dodge the
event, peaks at \MacClkCrushTermPeakPct\% for the expiry on the event
date, and decays; the clock reading (blue) is flat by construction.
(b)~One expiry (day 14) re-read each morning: the calendar reading
ramps from \MacClkCrushRampStartPct\% to \MacClkCrushPeakPct\% as the
event's four extra days weigh on fewer and fewer remaining calendar
days, then gives them up overnight --- $-$\MacClkCrushDropBp\ vol bp
--- when the event resolves.  No price surprised anyone at any
point.}
\label{fig:clkcrush}
\end{figure}

Panel~(b) is the view a desk actually watches: one fixed expiry ---
day 14, chosen to contain the event --- re-read every morning as the
week passes.  On valuation day $s$ the remaining calendar is $14-s$
days while the event is still ahead, so the remaining variance days
are $(14-s)+\MacClkClockExtraDays$, and the calendar reading of the
unchanged $30\%$ clock volatility is
\[
\sigma_{\rm cal}(s)
=30\%\times\sqrt{\frac{18-s}{14-s}}\,.
\]
Both the numerator and the denominator lose one day per morning, so
the \emph{fixed} four-day premium weighs on a shrinking base and the
ratio climbs: $30\%\sqrt{18/14}=\MacClkCrushRampStartPct\%$ on day
zero, then day by day up the staircase of panel~(b), reaching
$30\%\sqrt{8/4}=30\%\sqrt{2}=\MacClkCrushPeakPct\%$ on the last
morning before the report --- four calendar days left, eight variance
days.  That evening the event resolves.  The next morning one more
night has passed, so three calendar days remain, and the sum in
\eqref{eq:clkclock} no longer contains the event: three variance days
over three calendar days, and the reading is exactly $30\%$.  The
crush: $-$\MacClkCrushDropBp\ vol bp
overnight (computed unrounded: the peak is $30\sqrt2=42.43\%$), the
entire climb surrendered at once, and at no point did any price do
anything unscheduled.  The hump is a property of the ruler, not of
the prices.

Be precise about ``unscheduled,'' because the word carries the whole
point.  Prices \emph{do} change across the event night: the option
sheds the event's variance as it is realized, exactly as it sheds an
ordinary day's time value.  That is scheduled decay --- known to the
dollar in advance, priced from the start --- not news.  What the
construction lacks, and what the clock deliberately does not model,
is the \emph{surprise}: the realized move itself.  On a real board
the report's outcome repositions the whole surface; the clock only
ever reprices the schedule.  And the clock is a forecast --- if the
market decides mid-week that this quarter's report is worth six
extra days rather than four, the calendar's $N_e$ must be
re-estimated, and the reading of \emph{that} change is a genuine
repricing, not an artifact.  What the ruler explains is the part
every calendar-clock chart gets wrong on schedule: the climb and the
crush.

\section{The curve between the quotes}
\label{sec:clkinterp}

We can now read quoted expiries on the right clock.  But a board
quotes a handful of maturities, and half the daily uses of a term
structure need values \emph{between} them: plotting a dense curve,
marking an unlisted maturity, feeding a model that wants variance at
arbitrary dates.  Between the quotes something must be assumed, and
this section shows that the assumption is the clock choice again ---
in the sharpest form the chapter has to offer.

The rule the clock recommends is simple.  Interpolate \emph{total
variance}, not volatility --- variance is the additive quantity, the
thing that accumulates day by day, while a volatility is a ratio and
ratios do not add.  And interpolate it \emph{linearly in $\tau$}:
between two quoted expiries, let $w$ grow at a constant rate per
variance day.  Constant accrual per unit of the market's clock is
exactly the steady-diffusion picture of \cref{sec:clkwalk}, restated
between quotes; the alternative --- linear in $t$, constant accrual
per calendar day --- is the same picture on the wrong clock, and we
are about to measure the difference.

Build the test so that the truth is known.  A synthetic name runs at
a flat $30\%$ clock volatility --- the previous section's name ---
with one event worth \MacClkClockExtraDays\ extra days at day
\MacClkInterpEventDay.  Three expiries are quoted: 7, 30 and 90 days.
The event falls \emph{between} the second and third quotes: no quoted
expiry sits near it on either side.  Generate the three quotes from
the construction ($w=(0.30)^2\tau$ at each), then hand them --- plus
the event calendar --- to the two interpolation rules and ask each
for the whole dense curve.

\Cref{fig:clkinterp}(a) shows both answers as volatility readings
against calendar days.  Both curves pass through the three quoted
dots --- at a quote there is nothing to interpolate, so the rules
cannot disagree there.  Between 7 and 30 days they also agree: no
event inside, so a variance day is a calendar day and the two rules
coincide.  Everything happens between 30 and 90.  The blue curve ---
linear in $\tau$ --- runs flat at $30\%$ to day \MacClkInterpEventDay,
jumps as the event enters the maturity, and decays along the hump
shape of \cref{sec:clkcrush}: it reproduces the generating truth
exactly (to \MacClkInterpExactBp\ vol bp --- exactness is by
construction here, since the quotes sit on a curve that is itself
linear in $\tau$; the general statement closes the section).
The orange dashed curve --- linear in $t$ --- takes the interval's
total variance budget, which is correct at both endpoints, and
spreads it \emph{evenly across all sixty calendar days}.  The event's
four days of variance are smeared over the whole bracket.

Where does the smear hurt, and by how much?  Work the two sides of
the event date; the reading at maturity $T$ (in days) is
$30\%\times\sqrt{x(T)/T}$ with $x(T)$ the interpolated variance-day
count.  First the bracket's budget.  By the clock
\eqref{eq:clkclock}, the 30-day quote carries $x(30)=30$ variance
days and the 90-day quote carries $x(90)=90+4=94$ --- the event is
inside it --- so the bracket holds $94-30=64$ variance days over
$60$ calendar days, and the calendar rule pays them out evenly:
$\tfrac{64}{60}$ variance days per calendar day.  Just \emph{before} day 45, the truth carries no event:
$x_{\rm true}=45$.  The calendar rule has been paying the event out
steadily since day 30, so it carries
$x_{\rm cal}=30+15\cdot\tfrac{64}{60}=46$ --- one phantom variance
day --- and reads
$30\%\sqrt{46/45}=30.33\%$: \MacClkInterpOverBp\ vol bp of phantom
volatility at a maturity that contains no event at all.  Just
\emph{after} day 45, the truth jumps to $x_{\rm true}=49$ while the
calendar rule still shows $x_{\rm cal}=46$, three days short, and
reads $30\%\sqrt{46/45}$ against a truth of $30\%\sqrt{49/45}$:
\MacClkInterpUnderBp\ vol bp \emph{too low} on the first maturities
that do contain the event.  Panel~(b) plots this error across the
whole axis: zero wherever a quote or an eventless bracket pins the
rules together, a sawtooth around the event --- climbing to
$+\MacClkInterpOverBp$, snapping to $-\MacClkInterpUnderBp$, decaying
back to zero at the 90-day quote.

\begin{figure}[htbp]
\centering
\paperfig{\textwidth}{fig_clk_interp.pdf}
\caption{Interpolate in the clock, not the calendar (three quotes,
flat $30\%$ clock truth, event of \MacClkClockExtraDays\ extra days at
day \MacClkInterpOverDay).  (a)~Linear-in-$\tau$ (blue) reproduces the
truth exactly; linear-in-$t$ (orange, dashed) smears the event across
the whole 30--90-day bracket.  Both pass through every quote.
(b)~The calendar rule's reading error: $+\MacClkInterpOverBp$ vol bp
of phantom volatility just before the event date,
$-\MacClkInterpUnderBp$ vol bp just past it.}
\label{fig:clkinterp}
\end{figure}

Both signs are money.  A desk asked to price an option expiring the
day before the event --- an expiry that dodges it --- off the
calendar-interpolated curve quotes essentially the full
\MacClkInterpOverBp\ bp of volatility that does not exist; asked for
one expiring the day after, the one instrument whose whole point is
carrying the event, it quotes nearly the full \MacClkInterpUnderBp\
bp too little.  The errors sit precisely at the maturities an event
trader cares about, they point in opposite directions across one
calendar day, and no quoted expiry ever reveals them --- the rules
agree wherever a price exists.  Interpolation error between quotes
is invisible until it is traded against.

The general statement: on a real board the truth between
quotes is not known, and linear-in-$\tau$ is itself a convention ---
the cleanest available, being the constant-accrual statement on the
market's own clock, but a convention nonetheless.  What it
guarantees, and what the calendar rule cannot, is that \emph{known}
events stop leaking into maturities that do not contain them.  The
guarantee is two lines of algebra: between two quotes, $w$ is linear
in $\tau$, and by \eqref{eq:clkclock} $\tau$ grows linearly in $t$
except at an event date, where it jumps by $N_e/365$ --- so,
composed, the dense variance curve is linear in calendar time within
each stretch, pays each event's variance as a jump exactly at its
date, and changes slope only at the quotes.  The remaining question --- where the calendar's dates
and sizes come from --- is the chapter's last construction, and its
most delicate.

\section{Reading the clock off prices}
\label{sec:clkinverse}

Everything so far consumed a calendar: someone typed in ``four extra
days on August~26,'' and the clock did the rest.  For scheduled
earnings that is the right workflow --- the dates are published, and
a desk should use them.  But the sizes $N_e$ are nobody's published
number, and the board's own prices carry an opinion about them: we
\emph{saw} that opinion in \cref{fig:clkboard}(b), where one interval
accrued \MacClkBoardHotOverLull$\times$ the variance per day of its
neighbour.  This section turns that observation into an estimator ---
read the calendar \emph{off} the prices --- and, just as importantly,
measures what such an estimator can and cannot deliver.

Set up the objects.  The data is the quoted ladder: expiries
$t_1<\cdots<t_n$ with at-the-money total variances $w_1,\dots,w_n$
(the same per-interval increments $\Delta w_i=w_i-w_{i-1}$,
$\Delta t_i=t_i-t_{i-1}$ as before, with $w_0=t_0=0$).  The unknowns
are candidate event sizes: one candidate per interval, $N_i\ge0$
extra days somewhere inside interval $i$ --- where exactly inside
turns out not to matter, and a proposition below says why.  Given a
candidate vector $N=(N_1,\dots,N_n)$, each interval's width on the
variance clock is
\[
\Delta\tau_i=\Delta t_i+\frac{N_i}{365},
\]
and the chapter's central diagnostic gets its symbol: the
\emph{forward variance} of interval $i$,
\[
\mathcal{F}_i=\frac{\Delta w_i}{\Delta\tau_i}
\]
--- variance accrued per variance-clock year between two consecutive
expiries.  (Calligraphic $\mathcal{F}$; the roman $F$ remains the
forward price of Chapter~6.)  Two structural facts, each one line of
algebra, shape everything.

\emph{First: the data half of $\mathcal{F}_i$ is untouchable.}  The
increments $\Delta w_i$ are fixed by prices
(\cref{prop:clkgauge}(i)), and the calendar widths $\Delta t_i$ by
the listing.  The candidate $N$ enters only through the denominators
--- an estimator of the clock adjusts how the accrual is
\emph{divided}, never how much accrual there \emph{is}.

\emph{Second: an event can only pull its own interval down.}
Interval $j$'s width $\Delta\tau_j$ contains only interval $j$'s own
candidate --- an event in interval $i$ shifts the running clock
$\tau$ of every later expiry equally at both endpoints, so every
other width is unchanged.  And $N_i$ can only \emph{increase}
$\Delta\tau_i$, hence only \emph{decrease} $\mathcal{F}_i$.  The
estimator can lower a spike toward its neighbours; it cannot raise a
valley.  Its entire vocabulary is ``this interval's days were worth
more than ordinary days.''

Now the estimator itself.  A genuine pre-event board leaves a
specific signature: the event interval accrues faster than
\emph{both} its neighbours, so the ladder rises into it and falls out
of it --- a \emph{peak}.  An ordinary upward-sloping term structure
has no peak; a smoothly falling one has none either; markets show
both shapes constantly, and no event caused them.  So read the peaks,
and only the peaks.  Write $\mathcal{F}_i^{\rm ref}$ for interval
$i$'s \emph{reference}: the higher of its two neighbours,
\begin{equation}
\mathcal{F}_i^{\rm ref}=\max\big(\mathcal{F}_{i-1},\mathcal{F}_{i+1}\big),
\label{eq:clkref}
\end{equation}
evaluated on the calendar ladder ($N=0$).  The higher neighbour,
deliberately: an event is only ever charged with the excess over the
higher side, so a slope is never mistaken for a peak.  The rule is
then one inequality and one division.
\begin{equation}
N_i=365\,\Delta t_i\Big(\frac{\mathcal{F}_i}{\mathcal{F}_i^{\rm ref}}-1\Big)
\quad\text{if }
\frac{\mathcal{F}_i}{\mathcal{F}_i^{\rm ref}}-1\ge0.03
\text{ and } N_i\ge\tfrac12\text{ day},
\qquad N_i=0\text{ otherwise.}
\label{eq:clkpeak}
\end{equation}
The size is exact, not a compromise: substituting \eqref{eq:clkpeak}
into $\Delta\tau_i=\Delta t_i+N_i/365$ gives
$\Delta w_i/\Delta\tau_i=\mathcal{F}_i^{\rm ref}$ --- the event
interval is brought onto its higher neighbour and not one day
further.  The two floors are the rule's only constants.  The first
says a peak must stand at least three percent above its reference,
the second that the event it implies must be worth at least half a
day; on a short interval the second binds, on a long one the first.
Below either floor the excess is read as the fit noise of two ATM
variances, and nothing is installed.  When one interval is clipped
its neighbour's reference can only fall, so the rule is run as a
sweep --- candidates in ascending order, repeated until a sweep
changes nothing --- and lands on the largest ladder with no material
peak left: the smallest events that remove every peak.  Two edges
need a convention.  The last interval of the ladder is never a
candidate; it is the reference the interval before it is judged
against.  The first interval has no left neighbour: its reference is
its right neighbour, except when the ladder keeps falling beyond it
($\mathcal{F}_2>\mathcal{F}_3$), in which case that decay is
continued one interval forward at its log-linear rate --- a front
that decays smoothly from the first interval on is read as no event.
A dip is never raised (\emph{second} fact above), a plateau of two
equally hot intervals has no peak and is invisible, and two adjacent
hot intervals of unequal height are read as one event, on the
taller.  There is no objective and nothing to tune beyond the two
floors.  A flatness objective with a charge on installed days ---
the natural alternative --- shrinks every size it reports, by the
charge over the objective's curvature, and given enough candidates
flattens genuine term structure along with the peaks; the rule of
\eqref{eq:clkpeak} can do neither, because it only ever lowers a peak
to its neighbour.

Before trusting any output, be precise about what the data can
identify at all.

\begin{proposition}[Dates inside an interval are not identified]
\label{prop:clkident}
Every quantity the pipeline consumes --- fits, readings,
interpolation, the rule \eqref{eq:clkpeak} --- depends on the
event calendar only through the clock values $\tau(t_1),\dots,
\tau(t_n)$ at the quoted expiries, hence only through the per-interval
totals $N_i=\sum_{t_e\in(t_{i-1},t_i]}N_e$.  Consequently two events
inside the same interval are indistinguishable from one event of
their combined size placed anywhere in that interval, and an event's
date becomes identified only when a listed expiry separates it from
its neighbours.
\end{proposition}

\begin{proof}
The clock \eqref{eq:clkclock} at a quoted expiry $t_i$ is
$365\,t_i+\sum_{t_e\le t_i}N_e$: it reads the calendar only through
cumulative sums up to quoted dates.  Two calendars with the same
per-interval totals give the same sums at every $t_i$, hence the same
$\tau(t_i)$, the same $\Delta\tau_i$, and the same value of every
quantity listed.  The date within an interval never enters.
\end{proof}

So the rule's placement of its event --- mid-interval, by
convention --- is exactly that: a convention, recorded in the
contract table, not an estimate.  When the report date is public,
move the event there; the \emph{size} is what the board can be asked
for, and by the proposition it does not change.

What can the size estimator deliver?  \Cref{fig:clkident} is the
planted-event study: a flat clock volatility, one event of known
size inside one interval, apply \eqref{eq:clkpeak}, compare.  The
arithmetic is short enough to do first.  On a flat-clock board every
interval accrues the same variance per calendar day, $\sigma^2$, and
an event of $N^\ast$ extra days inside an interval of
$d=365\,\Delta t_i$ calendar days raises that interval's calendar
reading to $\sigma^2(1+N^\ast/d)$ while its neighbours stay at
$\sigma^2$.  The excess over the reference is therefore
$N^\ast/d$ exactly, whatever $\sigma$, and \eqref{eq:clkpeak} returns
\[
\widehat N=N^\ast
\quad\text{if}\quad N^\ast\ge\max\!\big(\tfrac12,\;0.03\,d\big),
\qquad\widehat N=0\ \text{otherwise}.
\]
Panel~(a) plots recovered against planted size and confirms it.  On a
dense board --- expiries from a week to a year, listed like a liquid
single name's, the event inside a \MacClkIdentDenseWidthD-day
interval --- the curve is the diagonal from
\MacClkIdentFloorDenseD\ days upward, exactly zero below; at $40\%$
volatility and at index-like $20\%$ the two curves coincide to
\MacClkIdentVolDiffD\ days.  On a coarse quarterly board, the event
inside a \MacClkIdentQuartWidthD-day interval, the curve is zero until
the planted event reaches \MacClkIdentFloorQuartD\ days and the
diagonal after.  Above the floors the largest recovery error on any
board is \MacClkIdentMaxErrD\ days --- rounding.  Handed a perfectly
flat ladder, the rule installs exactly \MacClkIdentFlatDays\ days:
silence is its default.

Panel~(b) draws the floor itself, $\max(\tfrac12,0.03\,d)$, against
the width of the bracketing interval.  Half a day until the intervals
reach seventeen days, three percent of the width beyond.  A two-day
event --- a small single-name report --- sits above the floor on any
interval up to \MacClkIdentWallWidthD\ days and below it on a
quarterly listing: the wall is in the listing, not in the volatility.
This is the price of an estimator with no objective.  It has no
shrinkage and no volatility dependence; what it cannot do is see an
event that a coarse ladder dilutes below three percent of an
interval's variance, and on such a board the date must be typed in.

\begin{figure}[htbp]
\centering
\paperfig{\textwidth}{fig_clk_ident.pdf}
\caption{The planted-event study.  (a)~Recovered vs planted size on
a dense board (event inside a \MacClkIdentDenseWidthD-day interval)
at $40\%$ and $20\%$ volatility, and on a quarterly board
(\MacClkIdentQuartWidthD-day interval): exact from the floor upward
(\MacClkIdentFloorDenseD\ and \MacClkIdentFloorQuartD\ days), exactly
zero below, identical at both volatilities.  (b)~The floor
$\max(\tfrac12,0.03\,d)$ against the interval width $d$; a two-day
event clears it up to \MacClkIdentWallWidthD-day intervals.}
\label{fig:clkident}
\end{figure}

We now know the estimator's character.  It reads exactly the kinks
the clock was built to explain; it can only pull hot intervals down;
it reports sizes exactly and stays silent below a floor set by the
listing's spacing; and it cannot date an
event more finely than the listing brackets it.  What it has never
been told is \emph{why} an interval runs hot.  The last section runs
it on the real frozen boards and confronts that limit directly.

\section{The clock of the frozen board}
\label{sec:clkfrozen}

The estimator is measured; point it at the real data.  This section
runs \eqref{eq:clkpeak} on the two frozen ladders --- NVDA, where the
chapter began, and SPY --- and reads the results within the limits
\cref{sec:clkinverse} measured.

\Cref{fig:clkread}(a) is the run with a horizon.  Candidates are allowed
at the first five expiries --- through December 2026, a
\emph{horizon} choice recorded in the contract --- while the later
intervals carry no candidates and serve only as references
(they sit beyond the panel's right edge; panel~(b) shows them).  The orange steps are the calendar ladder of
\cref{fig:clkboard}(b); the blue steps are the same increments
divided by the installed clock.  The rule finds one peak, exactly
where the chapter has been pointing: the earnings-bearing interval
(shaded) accrues \MacClkBoardHotFwd\ per year against a higher
neighbour of \MacClkReadHeroEarnAfter\ --- the September-to-December
interval, not the lull before it --- and is clipped onto it by
\MacClkReadHeroEarnD\ extra days, the whole budget:
\MacClkReadHeroTotalD\ days in total.  (Units, once: both ladders are
variance per year of their own clock --- calendar year for orange,
installed variance year for blue.)  Everything else is left standing,
orange peeking out behind blue.  The lull at
\MacClkReadHeroLullLevel\ is a dip, and a dip cannot be raised.  The
two short-dated intervals at the left edge --- the board's hot
Monday-to-Friday opening --- are not peaks: the second sits below
the first, and the first is judged against the back's decay
continued forward, \MacClkReadHeroFrontRef, which its own
\MacClkReadHeroFrontLevel\ does not exceed.  The in-horizon spread
of the ladder tightens from \MacClkReadHeroSpreadBeforeBp\ to
\MacClkReadHeroSpreadAfterBp\ variance bp; what remains is the
term structure's own shape.

Now read panel~(a) in two directions.  What the rule has
genuinely done is recover, \emph{from prices alone}, that the window
between August~21 and September~18 carries four ordinary days
more variance than its neighbour --- with no access to a news
feed.  The public schedule says why: NVDA's quarterly report sits in
that window.  But the rule does not know that, and cannot.  It
reads peaks; earnings, a product launch, a macro date, or a
genuinely lumpy term structure all look identical through
\eqref{eq:clkpeak}.  The installed calendar is an \emph{estimate of
sizes at conventional dates}, to be reviewed against the public
schedule --- not an oracle that discovered earnings.

\begin{figure}[htbp]
\centering
\paperfig{\textwidth}{fig_clk_read.pdf}
\caption{Reading the frozen boards' clocks (forward variance per
calendar year, orange, vs per installed variance year, blue).
(a)~NVDA with candidates through year-end: one peak, the earnings
interval (shaded), clipped onto its higher neighbour by
\MacClkReadHeroEarnD\ extra days; nothing else moves; the
in-horizon spread tightens from \MacClkReadHeroSpreadBeforeBp\ to
\MacClkReadHeroSpreadAfterBp\ variance bp.  (b)~NVDA with candidates
everywhere: a second peak, the December-to-March quarter, is clipped
by \MacClkReadFullMarchD\ days; \MacClkReadFullTotalD\ days in all,
the ladder's spread \MacClkReadFullSpreadBeforeBp\ to
\MacClkReadFullSpreadAfterBp\ variance bp --- not flat.
(c)~SPY: no peak clears the floor; nothing installed; the rising
ladder is left standing.}
\label{fig:clkread}
\end{figure}

Panel~(b) shows what happens when that review is skipped.  Same
board, same rule, but candidates at every expiry and no horizon.
The rule finds a second peak: the December-to-March interval accrues
\MacClkReadFullMarchBefore\ per year against a higher neighbour of
\MacClkReadFullMarchAfter, and is clipped by \MacClkReadFullMarchD\
days --- an ordinary three-month stretch whose elevated accrual is
some mixture of two more earnings reports and genuine term-structure
shape.  \MacClkReadFullTotalD\ days in all; the ladder's spread
falls from \MacClkReadFullSpreadBeforeBp\ variance bp to
\MacClkReadFullSpreadAfterBp, and no further, because the two dips
and the rising tail are not peaks and the rule cannot flatten what
it cannot raise.  Nothing malfunctioned.  A quarter that stands
above both its neighbours \emph{is} a peak, and the rule has no way
to say ``this interval is hot for a reason the clock should not
absorb.''  The horizon of panel~(a), and the human review of what it
installs, are not decorations on the method --- they are the method.

Panel~(c) is the control experiment the chapter owes the reader:
the index board.  SPY's calendar ladder rises steadily --- the
familiar upward term structure; its largest decrease anywhere is
\MacClkReadSpyDecBp\ variance bp --- and the rule, run with no
horizon at all, installs nothing, leaving
the ladder's \MacClkReadSpySpreadBp-variance-bp spread standing.
The mechanism is the floor.  A rising ladder has no peak at all,
and this one has exactly one, the two-to-four-day interval, which
stands \MacClkReadSpyPeakExcessPct\% above its higher neighbour ---
under the three-percent floor, and worth a fortieth of a day on a
two-day interval.  A diversified index dilutes any
single name's report to a fraction of a day; the rule's silence on
SPY is the study's flat-ladder result repeated on real data.  The
single name keeps an event clock; the index keeps the calendar.

Three paragraphs of scope complete the chapter's subject.  First,
below one day.  This chapter's smallest unit of time is the day, and
that too is a convention: variance does not accrue evenly
\emph{within} a day either.  The classic evidence is the weekend ---
measured across decades, an exchange holiday adds far less variance
than a trading day, because most variance arrives while markets
trade \citep{FrenchRoll1986} --- and the same asymmetry separates a
session hour from an overnight gap.  The day-weight idea continues
below the day boundary without new structure: a day becomes a
profile (so much of its weight during the session, so much
overnight), and a scheduled morning print --- a CPI release at 8:30
--- is simply an event whose date $t_e$ is a fraction of a day,
entering \eqref{eq:clkclock} unchanged.  The reference implementation
carries such a sub-day clock as an option: each trading day's unit
weight is split between the session and the hours around it (the
default profile is flat, $6.5/24$ to the session), and a non-trading
day keeps a stated weight, one by default --- so the weekend
discount of \citet{FrenchRoll1986} is a dial there, not a default,
and on the default whole-day clock a Friday board with steady
variance per trading day shows the interval after a weekend as a
peak per calendar day.  This book stops at day granularity,
stated here once.

Second, what the clock does not model.  The clock is a deterministic
forecast of the schedule of variance; it prices the \emph{anticipated}
part of an event and is silent on the realized surprise, on random
event sizes, and on volatility that is itself dynamic --- the
stochastic-volatility world, where an event is a random burst rather
than a known dilation.  That is model territory, developed in
\citet{Bergomi2016}, and it is the road this book declined at the end
of \cref{sec:clkboard}: richer, costlier, and unnecessary for the
readings this chapter set out to repair.

Third, where the idea comes from.  Reading asset prices on a
non-calendar clock is one of the oldest ideas in the subject:
\citet{Clark1973} explained heavy-tailed daily returns by
subordinating a Brownian motion to a transaction-volume clock, and
\citet{AneGeman2000} showed returns recover normality when read on
such an activity clock --- the empirical cousins of
\cref{thm:clkdds}.  This chapter's contribution to the reader is the
desk-sized version: a clock parameterized in days, priced into
readings by \cref{prop:clkgauge}, and estimable from the board with
stated limits.

\subsection{The contract}
\label{sec:clkcontract}

The chapter closes as the others do: three columns of decreasing
strength --- proved, measured, chosen.

\begin{table}[htbp]
\centering
\small
\caption{The chapter's contract.}
\label{tab:clkcontract}
\begin{tabular}{p{0.31\textwidth}p{0.31\textwidth}p{0.30\textwidth}}
\toprule
\textbf{Proved} & \textbf{Checked} & \textbf{Convention}\\
\midrule
The law at expiry knows the variance budget, not its schedule ---
exactly, for independent normal days; for every continuous martingale
by \cref{thm:clkdds} (cited). &
The walk of \cref{fig:clkwalk}: envelope kink
\MacClkWalkEnvKinkPct\% on the calendar axis, one smooth square root
on the variance axis. &
The clock $\tau$ is a deterministic forecast of the schedule;
realized surprises and random event sizes are out of scope
(\citealp{Bergomi2016}).\\
\addlinespace
Relabeling time (\cref{prop:clkgauge}): prices, butterfly validity
and calendar order invariant; readings covariant by
$\sqrt{t/\tau}$, interval rates by $\Delta t/\Delta\tau$. &
The worked week (\cref{fig:clkcrush}): ramp
\MacClkCrushRampStartPct\% to \MacClkCrushPeakPct\%, overnight crush
$-$\MacClkCrushDropBp\ vol bp, no unscheduled price move; the NVDA
dip (\MacClkBoardDipBp\ vol bp) diagnosed as the dodging expiry. &
Event currency: extra equivalent days $N_e$ per event, 365 days per
year; normalization \eqref{eq:clknorm} off throughout (a stated
budget choice).\\
\addlinespace
Linear-in-$\tau$ interpolation bends the dense curve at the
calendar's dates and nowhere else (\cref{sec:clkinterp}). &
The smear (\cref{fig:clkinterp}): $+\MacClkInterpOverBp$ vol bp
phantom before the event, $-\MacClkInterpUnderBp$ past it, zero at
every quote. &
Total variance linear in $\tau$ between quotes --- the
constant-accrual convention on the market's clock.\\
\addlinespace
Dates inside an interval are not identified
(\cref{prop:clkident}); an event can only lower its own interval's
$\mathcal{F}_i$; the peak rule \eqref{eq:clkpeak} recovers a planted
size exactly above its floor $\max(\tfrac12,0.03\,d)$. &
The study (\cref{fig:clkident}): exact to \MacClkIdentMaxErrD\,d
above floors of \MacClkIdentFloorDenseD/\MacClkIdentFloorQuartD\,d
(dense/quarterly), at every volatility; flat ladder
$\to$ \MacClkIdentFlatDays\ days. &
Candidate placement mid-interval; the reference as the higher
neighbour, the two floors and the horizon
(appendix~\ref{app:clkprotocol}); a plateau of equal peaks is
invisible.\\
\addlinespace
--- (this row is empirical by design) &
The frozen boards (\cref{fig:clkread}): \MacClkReadHeroEarnD\ days
onto NVDA's earnings interval (spread
\MacClkReadHeroSpreadBeforeBp$\to$\MacClkReadHeroSpreadAfterBp\ var
bp); the no-horizon run's second peak, \MacClkReadFullMarchD\ days
on an ordinary quarter; SPY left untouched. &
The installed calendar is reviewed against the public schedule
before use; the rule measures how much, never why.\\
\bottomrule
\end{tabular}
\end{table}

Step back to the book's thesis, one last time for Part~II.
\emph{Validity} was \cref{prop:clkgauge}: a clock is a relabeling,
and relabelings can never touch prices, densities, or order --- the
entire construction is arbitrage-proof by arithmetic.
\emph{Information} was what the board genuinely pins down: total
variances at quoted expiries, hence per-interval accrual rates ---
and, through the inverse problem, event \emph{sizes} above a
floor set by the listing, but never dates inside an interval and never causes.
\emph{Convention} was everything else, each piece now on the table:
the day currency, the normalization choice, the interpolation rule,
the candidate placement, the rule's two floors and its horizon.  With this
chapter, Part~II's work is complete: the raw board has been turned
into forwards and discounts (Chapter~6), European prices
(Chapter~7), and now readings on the market's clock --- the observation
is fully prepared.  What Part~III asks is what none of this statics
can answer: when the spot moves, what happens to the smile?  The
surface gives every derivative in the strike direction and none in
the spot direction --- a missing derivative, and the next chapter is
about the conventions that supply it.


\begin{chapterappendices}
\section{Protocol: data, constructions, and the rule's constants}
\label{app:clkprotocol}

Every number and figure in this chapter regenerates from one
deterministic script reading the frozen snapshot
(\texttt{scripts/ch08/gen\_figures.py}); nothing is refetched, and no
smile is refitted.  This appendix records the constructions.

\paragraph{The ATM ladders (\cref{fig:clkboard,fig:clkread}).}
For each (ticker, expiry) node of the frozen snapshot, the
at-the-money total variance $w_i$ is read from the node's prepared
mid quotes by linear interpolation of $w=\sigma^2 t$ in log-moneyness
$k$ at $k=0$, using the node's stored resolved forward (Chapter~6's
object) to define $k$.  No smoothing, no fitting: two bracketing
quotes decide each ladder point.  The maturities $t_i$ are the
snapshot's calendar year fractions.  Both names quote the same eight
expiries --- 2026-08-05, 08-07, 08-21, 09-18, 12-18, then 2027-03-19,
09-17 and 12-17: from the 2026-08-03 valuation date, 2, 4, 18, 46,
137, 228, 410 and 501 calendar days.  The NVDA 2-day node is the
snapshot's shortest and noisiest (Chapter~2's gallery); its ATM read
is used as stored, and none of the chapter's conclusions rests on it
alone.

\paragraph{The walk (\cref{fig:clkwalk}).}
\MacClkWalkDays\ trading days; ordinary daily return standard
deviation \MacClkWalkDailyPct\% ($\Delta w_i$ constant); the event
day (index \MacClkWalkEventDay) carries \MacClkWalkEventUnits\ ordinary
days' variance.  \MacClkWalkPaths\ paths from a fixed-seed generator
(seed 8) --- the chapter's only randomness; the envelopes are
computed, not simulated.

\paragraph{The worked calendar
(\cref{fig:clkclock,fig:clkcrush,fig:clkinterp}).}
One event, $N_e=\MacClkClockExtraDays$ extra days at day
\MacClkClockEventDay\ (\cref{fig:clkclock,fig:clkcrush}) or day
\MacClkInterpOverDay\ (\cref{fig:clkinterp}); flat clock volatility
$30\%$ where prices are needed; quoted expiries weekly to eight weeks
(\cref{fig:clkcrush}) or at $\{7,30,90\}$ days
(\cref{fig:clkinterp}).  The interpolation exhibit's quotes sit
exactly on $w=\sigma^2\tau$, so the linear-in-$\tau$ rule reproduces
the generator to \MacClkInterpExactBp\ vol bp by construction; the
figure's content is the calendar rule's deviation from it.

\paragraph{The peak rule (\cref{eq:clkpeak};
\cref{fig:clkident,fig:clkread}).}
Floors $0.5$ day and $3\%$ of the interval's variance; one
candidate per interval, dated mid-interval
(\cref{prop:clkident} makes the date a convention); the ladder's last
interval never a candidate; the first interval's reference is its
right neighbour, continued one interval forward at the log-linear
rate of the back's decay (measured in interval midpoints) when
$\mathcal{F}_2>\mathcal{F}_3>0$; candidates swept in ascending order
against the current ladder, each clipped to its reference when both
floors hold, the sweep repeated until nothing changes; the relative
excess is always that of the calendar ladder.  This is the reference
implementation's estimator verbatim, with one difference of input:
it reads each $w_i$ from the calibrated slice at $k=0$, the chapter
from the prepared quotes as above.  The study boards of
\cref{fig:clkident}: the dense board's expiries
are $\{7,14,28,56,91,182,365\}$ days with the planted event at day
40, the quarterly board's $\{91,182,273,365\}$ days with the event
at day 120.  The NVDA hero run of \cref{fig:clkread}(a) allows
candidates at expiries through 2026-12-18 (the horizon); the
panel~(b) and SPY runs allow candidates at every expiry.  Spreads
are max-minus-min of the forward-variance ladder, quoted in variance
bp; the in-horizon spread uses the intervals up to the horizon.

\paragraph{Units.}
Throughout the chapter: one variance bp is $10^{-4}$ of total
variance (or of a forward-variance rate per year); one vol bp is
$10^{-4}$ of absolute volatility; days convert to years at
$365$ calendar days per year, the same convention as the clock
\eqref{eq:clkclock}.

\end{chapterappendices}
